% ============================================================================
% 致谢
% ============================================================================

\section*{致谢}
\addcontentsline{toc}{section}{致谢}

本论文从选题、系统设计、模块实现到最终成文，得益于多方面的支持，在此一并致以诚挚的谢意。

首先感谢我的指导教师{{指导教师}}老师。从选题方向、系统抽象规范到论文叙述结构的耐心修订，{{指导教师}}老师都给予了细致而严谨的指导；对学术规范的坚持与对工程严谨性的要求，让我在反复推翻已写就的章节时仍保持耐心，在迷于「数据是否够漂亮」时回到「论证是否够诚实」的本心。

感谢{{所在学院}}各位任课教师，分布式系统、操作系统、软件工程与编程语言等课程打下的基础，使本文涉及的并发模型、事件驱动架构与隔离机制成为可能。

感谢开源社区。Rust 与 Tokio 提供的 \texttt{async/.await} 运行时与 \texttt{broadcast}/\texttt{mpsc} 通道、Axum 与 reqwest 的 HTTP 工具栈、SQLite 的 WAL 模式、Git 的 worktree 机制、ratatui 的 TUI 渲染框架，构成了 Fuxi 得以在本科阶段完成的根基。同样感谢 Anthropic Claude Code 与 OpenAI Codex 团队提供的 Agent CLI——它们既是 Fuxi 的接入对象，也是设计中反复参照的工程范本；感谢 A2A 协议工作组与 Model Context Protocol 项目发布的开放规范，是本论文核心贡献得以立足的语义基础。

需要在此坦诚说明 AI 协作在本工作中的角色。Fuxi 的相当一部分代码、测试、文档与本论文初稿，都是在与 Claude Code 等 AI 编程助手的密集协作下完成的，这一协作模式本身也是本论文的研究对象之一。但所有架构决策、协议设计、关键调试经验、性能取舍判断与论文叙述的最终结构均由本人独立做出并对其完整性负责——AI 助手的角色更接近「不知疲倦的结对编程伙伴」，不能替代研究者对系统目标与论证逻辑的判断。这一边界已在诚信声明中正式记录。

感谢同窗与朋友在系统压力测试、可用性反馈与论文初稿审阅中提出的真诚意见。最后感谢我的家人，是他们的长期支持让我能够专注于研究与写作。本论文的所有不足由本人独立承担。

\vspace{2cm}

\begin{flushright}
{\CJKfamily{heiti} 以琳}

2026 年 5 月于{{所在学院}}
\end{flushright}
